\section*{摘要}

尽管神经网络取得了显著成功，但其内部运行机制仍缺乏清晰的科学解释，黑箱问题依然存在。九项 nanoGPT 实验建立了训练—学习—推理（TLI）动力学关系：真实更新引起依赖接收状态的有限幅度功能响应；学习表现为分布式功能支撑的持久重组；目标跨越其特异性读出边界时，才会发生可观察的能力状态转变；推理则在不改变已学习状态的条件下调用这些支撑。第五项实验识别出的因素导出了一种直接预测更新后能力转变的算法，在 5,088 个目标上达到 91.43\% 的准确率和 91.49\% 的四类别宏平均召回率。随后，我们在另外五种神经网络架构中检验了这些关系及其预测。在第七种系统中，集中分配正确反馈削弱了其他已经掌握的技能；重新均衡反馈使这些技能的准确率从 70.31\% 恢复至 99.48\%。将推理产生的后果与后续训练连接起来，便形成统一的训练—学习—推理—强化（TLIR）闭环动力学。

\section{引言}

训练、泛化、能力形成和推理是神经网络研究的核心问题。不同研究领域已经发展出具有重要影响力的理论来研究这些问题。

在\textbf{优化理论}中，神经切线核（Neural Tangent Kernel，NTK）通过网络雅可比矩阵及相应的核，描述特定训练区域内参数变化引起的输出变化 \cite{jacot2018}。边缘稳定性（Edge of Stability）研究训练行为与 Hessian 最大特征值、学习率及损失轨迹之间的关系 \cite{cohen2021}。锐度感知最小化（Sharpness-Aware Minimization，SAM）利用参数邻域内的损失变化构造训练更新，以改善泛化性能 \cite{foret2021}。在这些研究中，\emph{当前状态}由参数及其局部导数或曲率表示；\emph{真实更新}由优化器产生；\emph{响应}则表现为模型输出或损失的变化。

在\textbf{表征学习与泛化}中，频谱偏置（Spectral Bias）考察函数或核模态随频率而异的学习速度 \cite{rahaman2019}。隐式偏置（Implicit Bias）研究优化器为何会选择特定的解，例如可分数据上的最大间隔方向 \cite{soudry2018}。神经坍缩（Neural Collapse）测量类别特征均值、类内变异性和分类器几何结构，揭示训练准确率达到饱和后表征仍可继续变化 \cite{papyan2020}。信息瓶颈（Information Bottleneck）理论研究输入、隐藏表征和目标之间的互信息，包括任务相关信息的保留与压缩 \cite{tishby2015}。这些研究中的\emph{内部组织}表现为频谱模态、优化器选择的解、特征几何或信息结构；\emph{功能变化}则表现为学习速度、解的演化和预测表现的变化。

在\textbf{能力形成}研究中，Grokking 描述了拟合训练样本与取得较强测试泛化能力之间的长时间延迟 \cite{power2022}。后续计算回路研究追踪了这一延迟期间内部计算结构的形成和重组 \cite{nanda2023}。双重下降（Double Descent）研究测试误差如何在接近插值区域时上升，并随模型规模、训练时间或数据条件变化而再次下降 \cite{belkin2019,nakkiran2020}。这些研究利用训练阶段或插值区域表征\emph{当前状态}，利用计算回路与解的演化表征\emph{内部变化}，并通过准确率和泛化误差衡量\emph{可观察能力}。

在\textbf{持续学习}中，灾难性遗忘是指新任务训练降低了先前任务的表现。弹性权重巩固（Elastic Weight Consolidation，EWC）使用先前任务的 Fisher 信息估计参数重要性，并惩罚重要参数的改变 \cite{kirkpatrick2017}。梯度情景记忆（Gradient Episodic Memory，GEM）使用新任务梯度与先前任务梯度的内积约束新任务更新 \cite{lopezpaz2017}。正交梯度下降（Orthogonal Gradient Descent，OGD）则将新梯度投影到能够近似保持先前任务输出的方向 \cite{farajtabar2020}。在这些方法中，\emph{已有状态}包含之前形成的参数和对函数变化敏感的方向；\emph{更新}是受到约束的新任务动作；\emph{相互作用}通过重叠的梯度或功能子空间产生；\emph{能力变化}通过保持、遗忘或迁移进行测量。

在\textbf{推理研究}中，注意力机制（Attention）利用查询（Query）、键（Key）和值（Value）计算依赖输入的信息权重 \cite{vaswani2017}。上下文学习（In-context Learning）研究在模型参数固定时，上下文中的示例和任务信息如何影响预测 \cite{brown2020}。思维链提示（Chain-of-Thought Prompting）研究中间推理步骤如何影响复杂任务表现 \cite{wei2022}。\emph{当前状态}是已训练模型；\emph{输入条件}包括查询、上下文和中间计算；\emph{内部计算}随注意力和激活状态改变；\emph{功能表现}则通过预测结果观察。

在\textbf{强化与反馈驱动训练}中，基于人类反馈的强化学习（Reinforcement Learning from Human Feedback，RLHF）利用偏好和奖励信号调整模型行为 \cite{ouyang2022}。自训练（Self-training）进一步使用模型预测生成训练数据 \cite{xie2020}。递归模型坍塌研究考察连续几代模型使用前代生成的数据进行训练时，低概率模态逐渐丧失的现象 \cite{shumailov2024}。这里，\emph{反馈}由奖励、偏好或生成数据分布提供；\emph{训练动作}是由此产生的更新；随后的\emph{功能变化}通过行为、准确率和生成分布的变化加以观察。

上述研究使用了不同的数学对象——核、Hessian 矩阵、梯度、类别特征、互信息、任务得分、注意力权重与奖励信号——但仍可识别出六类反复出现的共同变量：\emph{当前状态、真实更新、功能响应、内部组织、可观察能力和反馈}。\emph{TLIR 的核心关系是在开展上述跨领域比较之前，通过实验建立的；这里识别这些共同变量，是为了帮助读者理解理论，而不是描述理论最初的发现过程。}

图~\ref{fig:chain} 展示了相互连接的四种动力学角色。

\section{TLI 开环动力学的实验建立}

用于建立 TLI 的九项 nanoGPT 实验，所包含的原生干预、受控对照和独立验证，远比正文所能完整呈现的范围更广。本节概述这些实验的主要发现及由此建立的理论关系。完整实验协议、执行记录、负面结果和独立验证程序均保存在配套仓库中。

\subsection{实验仪器与证据控制}

这些实验所需的仪器，是随着具体测量问题的出现而逐步开发的，并非在理论建立之前就预先假定其存在。

这些仪器共同依托生成事实图（Generation-Fact Graph，GFG），它把记录的模型与优化器状态、真实更新事件、干预、输出及其来源身份连接起来。对于\textbf{训练}，真实更新因果分叉仪器恢复匹配的接收状态，对比完整更新、跳过更新、仅参数更新、优化器记忆变化和双向移植分支，从而区分更新本身与接收该更新的状态的作用。随后，有限幅度响应路径仪器将同一条已记录的真实更新，以预先冻结的不同幅度施加于模型，测量非线性数值响应和离散状态转变。对于\textbf{学习}，CSRG-4C 通过重复基线、单组件门控和双组件门控，测量已注册目标组功能支撑的必要性、备份、替代、协同作用及失效余量。最后，身份对齐目标边界账本把每个注册目标与精确的更新前状态和真实更新相对应，并将响应发生前的记录与更新后正确性的裁决分离。冻结协议、负面对照和独立验证程序进一步约束了对实验结果的解释。


表~\ref{tab:equation_guide} 概述十个编号公式的科学用途和主要符号含义。

\subsection{训练：真实更新与接收状态}

在一个已注册的训练步骤中，参数—优化器状态为 $Z_t=(\theta_t,\omega_t)$，真实参数更新定义为 $\Delta\theta_t=\theta_{t+1}-\theta_t$。

TL-E01 揭示了能力形成、下降和恢复现象，而这些现象不能仅通过损失、准确率或训练步数得到完整刻画。TL-E02 进一步表明，暂停有效的参数—优化器演化会延迟能力形成，而梯度裁剪阈值的对照实验并未复现这一效应。

TL-E03 使用完整更新/跳过更新分支及双向更新移植，建立了如下关系：相同真实更新作用于不同接收状态时，会产生不同的功能响应：
\begin{equation}
 R_q(\theta_A,\Delta\theta_B)\quad\text{vs.}\quad R_q(\theta_B,\Delta\theta_B),
 \label{eq:receiver}
\end{equation}
其中 $R_q(\theta,\Delta)=m_q(\theta+\Delta)-m_q(\theta)$。另行开展的优化器记忆交换检验了历史优化器状态如何影响更新的形成：
\begin{equation}
 \Delta^{(a,b)}=\operatorname{OptStep}(\theta_t,g_t,\omega_{a,b}).
 \label{eq:optmem}
\end{equation}
TL-E04 测量了一次真实更新所产生的完整有限幅度响应：
\begin{equation}
 r_{t,q}(\alpha)=m_q(\theta_t+\alpha\Delta\theta_t)-m_q(\theta_t).
 \label{eq:response}
\end{equation}
实验识别出近线性、饱和、加速、回转和符号反转等响应形式。随后，TL-E05 识别出三类主要响应条件因素：目标边界状态（F1）、目标特异性更新几何（F3）和参数—优化器接收状态（F5）。

\subsection{学习：分布式功能支撑与能力边界}

TL-E06 通过可逆的组件干预，建立了分布式功能支撑的实验证据，包括组件必要性、备份、替代、协同作用和失效余量。

随后，TL-E07 检验了真实训练更新如何持续重组这些支撑：
\begin{equation}
 \Delta P_{t,q}=P_q(\theta_t+\Delta\theta_t)-P_q(\theta_t).
 \label{eq:realloc}
\end{equation}
在 72 个实验区段中，记录了 1,636 次有效的支撑分配转变和 228 次主要支撑组件交接。

TL-E08 通过身份对齐的目标读出边界，将这些内部变化与可观察的能力联系起来。更新前后间隔（margin）的符号用于区分能力保持、丧失、未获得及获得或恢复。因此，学习被界定为\textbf{持久的功能支撑重组}，而能力转变由目标特异性读出边界的跨越决定。

\subsection{推理：冻结状态下对已学习支撑的调用}

INF-E01 研究了 13 条 nanoGPT 训练历史和 52 个检查点。在持久学习状态冻结的条件下，实验测量了依赖查询条件的组件效应、非加性的支撑交互，以及将已训练组件替换为其能力形成前的精确历史版本所产生的变化。

在固定的学习状态 $L^*$ 下，推理调用并组合已经学习形成的支撑：
\begin{equation}
 S_{q,c}^{\mathrm{active}}=\Pi_{q,c}(L^*),\qquad
 y_{q,c}=H_{q,c}(S_{q,c}^{\mathrm{active}}).
 \label{eq:inference}
\end{equation}
组件之间的非加性交互通过下式测量：
\begin{equation}
 I_{ij,q}=f_{\mathrm{full}}-f_{-i}-f_{-j}+f_{-i,-j},
 \label{eq:interaction}
\end{equation}
而对已学习组件版本的依赖通过下式检验：
\begin{equation}
 D_{j,q}=f(q;\theta^*)-f(q;\theta^{(j\leftarrow\mathrm{pre})}).
 \label{eq:rollback}
\end{equation}
精确恢复证明，这些干预并未持续改变已学习状态。实验由此建立了推理的动力学定义：\textbf{推理是在查询条件下调用并组合学习形成的功能支撑，而持久学习状态保持冻结。}

上述九项实验共同建立了相互连接的训练—学习—推理关系，形成 TLI 开环动力学。图~\ref{fig:tlir-evidence} 汇总了 TLIR 四种动力学角色的实验依据。
